\section*{Chapter \ref{ch:hf:hedging-and-inventory-in-practice} --- Hedging and Inventory in Practice}

\begin{solution}{exo:hf:hedging-and-inventory-in-practice:1}
$D=1\,000\times50\times1.2-500\times100\times0.8=60\,000-40\,000=\$20\,000$. It is less than half a contract (\$25\,000): no future is sold.
\end{solution}

\begin{solution}{exo:hf:hedging-and-inventory-in-practice:2}
$c_F/(2\lambda\sigma_F^2)=0.00005/(2\times10^{-4}\times10^{-4})=\$2\,500$; the hedge is $833\,747-2\,500=\$831\,247$.
\end{solution}

\begin{solution}{exo:hf:hedging-and-inventory-in-practice:3}
$2\times0.00005\times833\,747\approx\$83$.
\end{solution}

\begin{solution}{exo:hf:hedging-and-inventory-in-practice:4}
The idiosyncratic gaps: each name moves by its own news overnight (1\% here, fat-tailed), and $\sqrt{\sum_i(q_iS_i\sigma_i)^2}$ over forty positions
of about \$40\,000 is about \$3\,000. Only trading the stocks removes it.
\end{solution}

\begin{solution}{exo:hf:hedging-and-inventory-in-practice:5}
$b^\ast=\bigl(3c\sigma_D^2/(4\lambda\sigma_F^2)\bigr)^{1/3}$: \$613\,000, \$425\,000 and \$284\,000. A higher charge on variance makes unhedged exposure
dearer relative to the cost of trading, so the desk tolerates less of it before hedging; the cube root makes the band change slowly.
\end{solution}

\begin{solution}{exo:hf:hedging-and-inventory-in-practice:6}
The desk still earns its half-spread on every fill; it only tilts which side trades, and clients on the favoured side cross a quote that is a little
better for them. It is not free when the flow is one-sided (everyone sells at 15:30), when the skew moves the quote through the far touch and the
desk pays the spread, or when the book is too large for the remaining hour's flow.
\end{solution}

\begin{solution}{exo:hf:hedging-and-inventory-in-practice:7}
With the future at one basis point a dollar the formula's band is \$358\,000 at $\lambda=10^{-4}$; among the simulated bands, \$500\,000 scores best
(\$41\,540, against \$40\,680 at \$250\,000 and \$40\,490 at \$1 million), and hedging at every move falls to \$31\,560. The optimum moves out by about the
cube root of two.
\end{solution}

\begin{solution}{exo:hf:hedging-and-inventory-in-practice:8}
Minimising risk is not the objective; risk has a price and so does hedging. The rule trades 2\,682 contracts a day for \$6\,700, 14\% of the spread
income, and removes only slightly more risk than a \$250\,000 band that costs \$1\,880. It also leaves the idiosyncratic risk untouched.
\end{solution}

\begin{solution}{pb:hf:hedging-and-inventory-in-practice:1}
\begin{enumerate}
\item See \cref{def:hf:hedging-and-inventory-in-practice:instrument}: the cost of trading it and the share of the risk it removes.
\item See \cref{def:hf:hedging-and-inventory-in-practice:delta}: $D=\sum_iq_iS_i\beta_i$.
\item $D^2\sigma_F^2+\sum_i(q_iS_i\sigma_i)^2$.
\item After a fall, fills in every name tend to make the desk long at once: the inventories share the factor's sign.
\item See \cref{prop:hf:hedging-and-inventory-in-practice:dead} and its proof.
\item $b^\ast=\bigl(3c\sigma_D^2/(4\lambda\sigma_F^2)\bigr)^{1/3}$, \$284\,000.
\item Basis between the future and the index, betas estimated with error, and stocks outside the index.
\item Offsetting positions cancel before anyone pays a spread; the books must be summed in near real time and each residual must stay within its own
limits.
\item None: \$0 and \$8\,540 a day; future at every move: \$6\,700 and \$4\,410; band \$250\,000: \$1\,880 and \$4\,760; stocks: \$12\,760 and \$2\,525.
\item The \$250\,000 band (\$42\,570); the formula gives \$284\,000, between the two best simulated bands.
\item Price pressures of 49 basis points with a half life of 0.92 days on average (Hendershott and Menkveld); spreads wider when specialists hold
large positions or lose money (Comerton-Forde and co-authors).
\item \euro1.55 a trade on the spread net of fees; \euro0.68 lost on positions.
\item Flatten: \$438 and no risk; future overnight: \$83, standard deviation \$3\,086, 99\% loss \$7\,358, expected shortfall \$8\,813; keep: 99\% loss
\$14\,776, expected shortfall \$19\,795, standard deviation \$5\,879.
\item Flattening (\$438 against \$1\,035 and \$3\,454): the idiosyncratic risk left by the future costs more in a night than the flattening does.
\item See \cref{def:hf:hedging-and-inventory-in-practice:flatten}: through the quotes, the closing auction, the continuous market, the future
overnight.
\item It halves the book at the close and the bill, from \$381 to \$163 on average.
\item When the book is small relative to the cost of flattening, when tomorrow's flow is expected to unwind it, or when the positions are a hedge of
something else the firm holds.
\item The future hedge leaves a factor residual proportional to the beta errors; the choice tilts towards stocks or towards a smaller hedge.
\item Each option's delta times its underlying's beta and price; gamma makes the number move with the market.
\item When the risk it removes is worth more, at the firm's price of risk, than the spread and fees it pays.
\end{enumerate}
\end{solution}

\begin{solution}{iq:hf:hedging-and-inventory-in-practice:1}
Sell in the closing auction, sell the stocks now, sell a bank ETF or the index future against them, or keep. The sector ETF removes more of the
sector risk than the index future; the closing auction avoids the spread but is a single price that a large order moves. The answer depends on the
size relative to the names' closing volume and on the desk's limits.
\iqlookfor{routes, costs, residual risk, size.}
\end{solution}

\begin{solution}{iq:hf:hedging-and-inventory-in-practice:2}
Regress each name's returns on the future's over a window long enough to be stable and short enough to be current (with care for asynchronous
closes); test by the variance of the hedged book out of sample, compared with the unhedged.
\iqlookfor{estimation and an out-of-sample test.}
\end{solution}

\begin{solution}{iq:hf:hedging-and-inventory-in-practice:3}
With a proportional cost, a small trade costs little, so the desk trades just enough to stay inside the band; with a fixed cost, each trade costs
the same whatever its size, so when the desk trades it goes all the way to the best position.
\iqlookfor{marginal versus fixed cost.}
\end{solution}

\begin{solution}{iq:hf:hedging-and-inventory-in-practice:4}
The three exposures may have the same sign and add up to a firm-wide exposure beyond any limit; or they may offset each other, and each desk hedges
anyway, paying for three hedges of nothing. The firm needs the sum in real time.
\iqlookfor{aggregation across desks.}
\end{solution}

\begin{solution}{iq:hf:hedging-and-inventory-in-practice:5}
Positions from every book's fills (drop copies), a shared table of betas and deltas updated on a timer, prices from the fair-value service, a sum
updated per fill; publish to the hedger and risk; reconcile against the end-of-day positions.
\iqlookfor{event-driven aggregation with reference data.}
\end{solution}

\begin{solution}{iq:hf:hedging-and-inventory-in-practice:6}
Minimise $c|x|+\lambda h\sigma_F^2(D-x)^2$: the half-width is $c/(2\lambda h\sigma_F^2)$, shrinking as the horizon $h$ grows, since the variance
avoided grows with the time the exposure is held while the cost of the trade does not.
\iqlookfor{the first-order condition and the horizon.}
\end{solution}
